% problem_id: S001-theorem-inifinite-galois-theory
% source_id: S001 (Stacks Project)
% source_file: fields.tex
% statement_locator: fields.tex lines 3053--3068
% proof_locator: fields.tex lines 3070--3134
% env: theorem
% id_note: label 在本来源内唯一
% status: complete (statement + author proof verbatim + reason + locators)
%
\begin{theorem}[Fundamental theorem of infinite Galois theory]
\label{theorem-inifinite-galois-theory}
Let $L/K$ be a Galois extension. Let $G = \text{Gal}(L/K)$
be the Galois group viewed as a profinite topological group
(Lemma \ref{lemma-galois-profinite}). Then we have $K = L^G$ and the map
$$
\{\text{closed subgroups of }G\}
\longrightarrow
\{\text{subextensions }L/M/K\},\quad
H \longmapsto L^H
$$
is a bijection whose inverse maps $M$ to $\text{Gal}(L/M)$.
The finite subextensions $M$ correspond exactly to the open
subgroups $H \subset G$. The normal closed subgroups $H$ of $G$
correspond exactly to subextensions $M$ Galois over $K$.
\end{theorem}

\begin{proof}
We will use the result of finite Galois theory
(Theorem \ref{theorem-galois-theory})
without further mention.
Let $S \subset L$ be a finite subset. There exists a tower
$L/E/K$ such that $K(S) \subset E$ and such that
$E/K$ is finite Galois, see Lemma \ref{lemma-normal-closure-inside-normal}.
In other words, we see that $L/K$ is the union of its finite
Galois subextensions.
For such an $E$, by Lemma \ref{lemma-galois-infinite}
the map $\text{Gal}(L/K) \to \text{Gal}(E/K)$ is surjective
and continuous, i.e., the kernel is open because the topology
on $\text{Gal}(E/K)$ is discrete.
In particular we see that no element of $L \setminus K$ is fixed by
$\text{Gal}(L/K)$ as $E^{\text{Gal}(E/K)} = K$.
This proves that $L^G = K$.

\medskip\noindent
By Lemma \ref{lemma-galois-goes-up} given a subextension $L/M/K$
the extension $L/M$ is Galois. It is immediate from the definition
of the topology on $G$ that the subgroup $\text{Gal}(L/M)$ is closed.
By the above applied to $L/M$ we see that $L^{\text{Gal}(L/M)} = M$.

\medskip\noindent
Conversely, let $H \subset G$ be a closed subgroup. We claim that
$H = \text{Gal}(L/L^H)$. The inclusion $H \subset \text{Gal}(L/L^H)$
is clear. Suppose that $g \in \text{Gal}(L/L^H)$. Let $S \subset L$
be a finite subset. We will show that the open neighbourhood
$U_S(g) = \{g' \in G \mid g'(s) = g(s)\}$ of $g$ meets $H$.
This implies that $g \in H$ because $H$ is closed.
Let $L/E/K$ be a finite Galois subextension containing $K(S)$
as in the first paragraph of the proof and consider the homomorphism
$c : \text{Gal}(L/K) \to \text{Gal}(E/K)$.
Then $L^H \cap E = E^{c(H)}$. Since $g$ fixes $L^H$ it fixes
$E^{c(H)}$ and hence $c(g) \in c(H)$ by finite Galois theory.
Pick $h \in H$ with $c(h) = c(g)$. Then $h \in U_S(g)$ as desired.

\medskip\noindent
At this point we have established the correspondence between closed
subgroups and subextensions.

\medskip\noindent
Assume $H \subset G$ is open. Arguing as above we find that
$H$ contains $\text{Gal}(L/E)$ for some large enough finite
Galois subextension $E$ and we find that $L^H$ is contained
in $E$ whence finite over $K$. Conversely, if $M$ is a finite
subextension, then $M$ is generated by a finite subset $S$
and the corresponding subgroup is the open subset $U_S(e)$
where $e \in G$ is the neutral element.

\medskip\noindent
Assume that $K \subset M \subset L$ with $M/K$ Galois.
By Lemma \ref{lemma-galois-infinite} there is a surjective
continuous homomorphism of Galois groups
$\text{Gal}(L/K) \to \text{Gal}(M/K)$ whose
kernel is $\text{Gal}(L/M)$. Thus $\text{Gal}(L/M)$ is a normal
closed subgroup.

\medskip\noindent
Finally, assume $N \subset G$ is normal and closed. For any
$L/E/K$ as in the first paragraph of the proof, the image
$c(N) \subset \text{Gal}(E/K)$ is a normal subgroup.
Hence $L^N = \bigcup E^{c(N)}$ is a union of Galois extensions
of $K$ (by finite Galois theory) whence Galois over $K$.
\end{proof}
